\subsection{Same Cardinality, Finite, Countable, and Uncountable}

\begin{definition} \label{def:cardinality-equivalence} \hlt{Same Cardinality}\\
For sets $A,B$, define $A\sim_{\#}B$ iff there exists a bijection $f:A\to B$.\\
Then $A$ and $B$ have \hlt{same cardinality}, or are \hlt{equivalent}. The relation $\sim_{\#}$ is an equivalence relation.\\
$\operatorname{id}_A$ proves reflexivity, inverses of bijections prove symmetry, and composition of bijections proves transitivity.
\end{definition}

\begin{definition} \label{def:finite-countable-uncountable} \hlt{Finite, Countable, and Uncountable}\\
Set $J_0\defeq\emptyset$, and for $n\geq1$ set $J_n\defeq\{1,2,\ldots,n\}$. A set $A$ is
\begin{enumerate}[label=(\roman*)]
\setlength{\itemsep}{0pt}
\item \hlt{finite} if $A\sim_{\#} J_n$ for some $n\in\N$;
\item \hlt{infinite} if it is not finite;
\item \hlt{countable} if $A\sim_{\#} \N$;
\item \hlt{uncountable} if it is neither finite nor countable;
\item \hlt{at most countable} if it is finite or countable.
\end{enumerate}
Countable sets are also called enumerable or denumerable.
\end{definition}

\begin{proposition} \label{prop:finite-cardinality-well-defined} \hlt{Finite Cardinality is Well Defined}\\
Let $m,n\in\N$. If $A\sim_{\#}J_m$ and $A\sim_{\#}J_n$, then $m=n$.\\
Hence for finite $A$, the value $\#A\defeq n$ with $A\sim_{\#}J_n$ is well-defined.
\begin{proof} By transitivity of $\sim_{\#}$, it suffices to prove $J_m\sim_{\#}J_n\Rightarrow m=n$.\\
Induct on $m$. If $m=0$, then $J_m=\emptyset$, so any bijection $J_m\to J_n$ forces $J_n=\emptyset$ and $n=0$.\\
Assume the claim for $m-1$, with $m\geq1$, and let $f:J_m\to J_n$ be bijective. Then $n\geq1$.\\
Removing $m$ and $f(m)$ gives a bijection $J_m\setminus\{m\}\to J_n\setminus\{f(m)\}$. Since these sets have the standard relabelings $J_{m-1}$ and $J_{n-1}$, the induction hypothesis gives $m-1=n-1$, hence $m=n$.\\
\hlt{Structural Commentary}: Finite size is bijection-invariant.\\
\hlt{Key Insight}: Remove one matched element.
\end{proof}
\end{proposition}

\begin{definition} \label{def:finite-set-summation} \hlt{Summation over a Finite Set}\\
Let $X$ be a finite set with $\#X=n$, and let $f:X\to\C$.\\
Choose any bijection $g:J_n\to X$, where $J_n=\{1,\ldots,n\}$, and set $\sum\limits_{x\in X}f(x)\defeq\sum_{i=1}^{n}f(g(i))$.\\
When $X=\emptyset$, this is the empty sum $0$.\\
\hlt{Key Insight}: Proposition \ref{prop:finite-set-summation-laws} shows that the value is independent of the chosen bijection.
\end{definition}

\begin{proposition} \label{prop:finite-set-summation-laws} \hlt{Finite Set Summation Laws}\\
Let $X,Y$ be finite sets.
\begin{enumerate}[label=(\roman*)]
\setlength{\itemsep}{0pt}
\item \hlt{Well-definedness}: different bijections $J_{\#X}\to X$ give the same value in Definition \ref{def:finite-set-summation}.
\item \hlt{Substitution}: if $h:Y\to X$ is a bijection and $f:X\to\C$, then $\sum\limits_{x\in X}f(x)=\sum\limits_{y\in Y}f(h(y))$.
\item \hlt{Linearity}: for $f,g:X\to\C$ and $c\in\C$, $\sum\limits_{x\in X}(f(x)+g(x))=\sum\limits_{x\in X}f(x)+\sum\limits_{x\in X}g(x)$ and $\sum\limits_{x\in X}cf(x)=c\sum\limits_{x\in X}f(x)$.
\item \hlt{Disjoint unions}: if $X\cap Y=\emptyset$ and $f:X\cup Y\to\C$, then $\sum\limits_{z\in X\cup Y}f(z)=\sum\limits_{x\in X}f(x)+\sum\limits_{y\in Y}f(y)$.
\item \hlt{Finite Fubini}: if $f:X\times Y\to\C$, then $\sum\limits_{x\in X}\sum\limits_{y\in Y}f(x,y)=\sum\limits_{(x,y)\in X\times Y}f(x,y)=\sum\limits_{y\in Y}\sum\limits_{x\in X}f(x,y)$.
\end{enumerate}
\begin{proof}
\begin{enumerate}[label=(\roman*)]
\setlength{\itemsep}{0pt}
\item For $f:X\to\C$, let $n=\#X$ and let $g,k:J_n\to X$ be bijections. Induct on $n$. The case $n=0$ is empty. If $n\geq1$, let $j\in J_n$ satisfy $k(j)=g(n)$. Removing $n$ from $g$ and $j$ from $k$ gives bijections onto $X\setminus\{g(n)\}$; after relabeling by $J_{n-1}$, the induction hypothesis makes the remaining sums equal. Adding the common last term $f(g(n))=f(k(j))$ gives equality of the two sums.
\item Let $g:J_{\#Y}\to Y$ be a bijection. Then $h\circ g:J_{\#Y}\to X$ is a bijection, so $\#Y=\#X$ by Proposition \ref{prop:finite-cardinality-well-defined}.\\
Hence, by (i) and Definition \ref{def:finite-set-summation}, $\sum\limits_{x\in X}f(x)=\sum\limits_{i=1}^{\#Y}f(h(g(i)))=\sum\limits_{y\in Y}f(h(y))$.
\item Choose a bijection $\gamma:J_{\#X}\to X$.\\
Proposition \ref{prop:finite-sum-basic-laws} gives $\sum\limits_{x\in X}(f(x)+g(x)) =\sum\limits_{x\in X}f(x)+\sum\limits_{x\in X}g(x)$ and $\sum\limits_{x\in X}cf(x)=c\sum\limits_{x\in X}f(x)$
\item Let $\alpha:J_{\#X}\to X$ and $\beta:J_{\#Y}\to Y$ be bijections, and enumerate $X\cup Y$ by listing $\alpha$ first and then $\beta$. Since $X\cap Y=\emptyset$, this is a bijection, and Proposition \ref{prop:finite-sum-basic-laws} gives $\sum\limits_{z\in X\cup Y}f(z)=\sum\limits_{x\in X}f(x)+\sum\limits_{y\in Y}f(y)$.
\item Let $\alpha:J_{\#X}\to X$ and $\beta:J_{\#Y}\to Y$ be bijections.\\
The maps $\ell((i-1)\#Y+j)=(\alpha(i),\beta(j))$ and $r((j-1)\#X+i)=(\alpha(i),\beta(j))$ for $i\in J_{\#X}$ and $j\in J_{\#Y}$ give the two lexicographic bijections $J_{(\#X)(\#Y)}\to X\times Y$.\\
Applying Definition \ref{def:finite-set-summation}, Proposition \ref{prop:finite-sum-basic-laws}, and (i) yields $\sum\limits_{x\in X}\sum\limits_{y\in Y}f(x,y) =\sum\limits_{(x,y)\in X\times Y}f(x,y) =\sum\limits_{y\in Y}\sum\limits_{x\in X}f(x,y)$.
\end{enumerate}
\hlt{Structural Commentary}: A finite set sum is an order-free finite sum.\\
\hlt{Key Insight}: Bijections are bookkeeping devices; the value does not depend on the bookkeeping order.
\end{proof}
\end{proposition}

\begin{definition} \label{def:sequence-as-enumeration} \hlt{Sequences as Enumerations}\\
Let $A$ be a set.
\begin{enumerate}[label=(\roman*)]
\setlength{\itemsep}{0pt}
\item A \hlt{sequence} in $A$ is a map $x:\N\to A$, written $(x_n)_{n\in\N}$ with $x_n\defeq x(n)$.
\item A sequence $x:\N\to A$ \hlt{enumerates} $A$ iff $x(\N)=A$.
\item A sequence $x:\N\to A$ is a \hlt{distinct enumeration} of $A$ iff $x:\N\to A$ is bijective.
\end{enumerate}
Thus $A$ is countable iff it has a distinct enumeration.
\end{definition}

\begin{example} \label{ex:integers-are-countable} \hlt{The Integers Are Countable}\\
The set $\Z$ is countable. One explicit enumeration is given by
\begin{align}
f(n)=
\begin{cases}
\frac{n}{2}, & n \text{ even},\\[2mm]
-\frac{n-1}{2}, & n \text{ odd},
\end{cases}
\qquad n\in\N. \nonumber
\end{align}
This lists $\Z$ as $0,1,-1,2,-2,3,-3,\ldots$.
\end{example}

\begin{theorem} \label{thm:infinite-subset-of-countable-is-countable} \hlt{Infinite Subsets of Countable Sets}\\
Let $A$ be countable and let $E\subseteq A$ be infinite. Then $E$ is countable.
\begin{proof} Let $x:\N\to A$ be a distinct enumeration of $A$. Since $E$ is infinite, $I_E\defeq\{n\in\N:x_n\in E\}$ is infinite.\\
Write $I_E=\{n_1<n_2<\cdots\}$ and define $g:\N\to E$ by $g(k)\defeq x_{n_k}$. The terms $(x_n)$ are distinct, so $g$ is injective.\\
If $y\in E$, then $y=x_n$ for a unique $n\in\N$; since $n\in I_E$, $n=n_k$ for some $k$, and $y=g(k)$.\\
Thus $g$ is bijective, so $E$ is countable.\\
\hlt{Structural Commentary}: Subsets select indices.\\
\hlt{Key Insight}: Keep the hitting indices.
\end{proof}
\end{theorem}

\begin{proposition} \label{prop:subsets-and-images-at-most-countable} \hlt{Subsets and Images of At Most Countable Sets}\\
Let $A$ be at most countable.
\begin{enumerate}[label=(\roman*)]
\setlength{\itemsep}{0pt}
\item If $B\subseteq A$, then $B$ is at most countable.
\item If $f:A\to Y$ is a function, then $f(A)$ is at most countable.
\end{enumerate}
\begin{proof}
If $A$ is finite, both claims are immediate. Otherwise write $A=\{a_n:n\in\N\}$ with the $a_n$ distinct.
\begin{enumerate}[label=(\roman*)]
\setlength{\itemsep}{0pt}
\item If $B$ is infinite, let $I_B\defeq\{n:a_n\in B\}=\{n_1<n_2<\cdots\}$.\\
Then $k\mapsto a_{n_k}$ bijects $\N$ onto $B$; the finite case is trivial.
\item Let $I\defeq\{n:f(a_n)\notin\{f(a_m):m<n\}\}$. By (i), $I$ is at most countable, $n\mapsto f(a_n)$ bijects $I$ onto $f(A)$.
\end{enumerate}
\hlt{Structural Commentary}: Subsets select indices; images select first appearances.\\
\hlt{Key Insight}: Subsets and images preserve at-most-countability.
\end{proof}
\end{proposition}

\begin{proposition} \label{prop:at-most-countable-iff-enumerable-image} \hlt{At Most Countable Sets Are Exactly Enumerable Images}\\
For a set $B$, the following are equivalent:
\begin{enumerate}[label=(\roman*)]
\setlength{\itemsep}{0pt}
\item $B$ is at most countable.
\item $B=\emptyset$ or there exists a surjection $f:\N\to B$.
\end{enumerate}
Thus a nonempty set is at most countable iff it is the image of a sequence.
\begin{proof}
\begin{enumerate}[label=(\roman*)]
\setlength{\itemsep}{0pt}
\item[$(i)$]$\Rightarrow(ii)$ If $B=\emptyset$, (ii) holds. If $B$ is countable, choose a bijection $f:\N\to B$. If $B=\{b_1,\ldots,b_m\}$ is finite and nonempty, define $f(0)=b_1$, $f(n)=b_n$ for $1\leq n\leq m$, and $f(n)=b_m$ for $n>m$.\\
In both nonempty cases, $f:\N\to B$ is surjective.
\item[$(ii)$]$\Rightarrow(i)$ If $B=\emptyset$, then $B$ is finite.\\
If $f:\N\to B$ is surjective, then $B=f(\N)$ is at most countable by Proposition \ref{prop:subsets-and-images-at-most-countable}.
\end{enumerate}
\hlt{Structural Commentary}: Surjections list with repetition.\\
\hlt{Key Insight}: At most countable means sequence image.
\end{proof}
\end{proposition}

\begin{figure}[H]
  \centering
  \begin{tikzpicture}[
  >=Stealth,
  every node/.style={font=\small},
  entry/.style={draw=gray!55, fill=gray!5, rounded corners=1pt, minimum width=0.62cm, minimum height=0.40cm},
  active/.style={draw=orange!60!black, fill=orange!10},
  scan/.style={->, blue!70!black, line width=0.7pt},
  guide/.style={orange!75!black, line width=0.65pt},
  output/.style={draw=gray!55, fill=gray!5, rounded corners=1pt, inner xsep=5pt, inner ysep=4pt}
]
  \draw[guide] (1.70,-0.65) -- (0.85,-1.30);
  \draw[guide] (2.55,-0.65) -- (0.85,-1.95);
  \draw[guide] (3.40,-0.65) -- (0.85,-2.60);

  \foreach \row/\name in {1/E_1,2/E_2,3/E_3,4/E_4} {
    \node[anchor=east] at (0,-0.65*\row) {$\name$};
    \foreach \col in {1,...,5} {
      \node[entry] (x-\row-\col) at (0.85*\col,-0.65*\row) {$x_{\row\col}$};
    }
    \node[anchor=west, gray!65] at (4.55,-0.65*\row) {$\cdots$};
  }

  \foreach \row/\col in {1/1,1/2,2/1,1/3,2/2,3/1,1/4,2/3,3/2,4/1} {
    \node[entry, active] at (0.85*\col,-0.65*\row) {$x_{\row\col}$};
  }

  \node[output, anchor=west] (list) at (5.25,-1.45)
    {$x_{11},\,x_{12},\,x_{21},\,x_{13},\,x_{22},\,x_{31},\,x_{14},\,\cdots$};
  \draw[scan] (4.55,-1.45) -- (list.west);
\end{tikzpicture}

  \caption{Diagonal scan of countable rows into one sequence, skipping repeats.}
  \label{fig:countable-union-diagonal-scan}
\end{figure}

\begin{theorem} \label{thm:countable-union-of-countable-sets}\label{thm:cardinality-main-results-source-todo} \hlt{Countable Union of Countable Sets}\\
Let $(E_n)_{n\geq1}$ be countable sets and set $S\defeq\bigcup\limits_{n=1}^{\infty}E_n$. Then $S$ is countable.
\begin{proof} For each $n\geq1$, choose a bijection $x_n:\N\to E_n$ and write $x_{n,k}\defeq x_n(k)$.\\
Let $D\defeq\{(n,k):n\geq1,\ k\in\N\}$. Diagonal order gives a surjection $d:\N\to D$.\\
Set $a(t)\defeq x_{n,k}$ for $d(t)=(n,k)$; then $a(\N)=S$, so $S$ is at most countable by Proposition \ref{prop:at-most-countable-iff-enumerable-image}.\\
Since $E_1\subseteq S$ and $E_1$ is countable, $S$ is infinite. Hence $S$ is countable.\\
\hlt{Structural Commentary}: Diagonal scanning enumerates pairs.\\
\hlt{Key Insight}: Scan the two-index list.
\end{proof}
\end{theorem}

\begin{corollary} \label{cor:at-most-countable-union-at-most-countable} \hlt{At Most Countable Unions}\\
If $A$ is at most countable and $B_\alpha$ is at most countable for each $\alpha\in A$, then $\bigcup\limits_{\alpha\in A}B_\alpha$ is at most countable.
\begin{proof} Let $T\defeq\bigcup\limits_{\alpha\in A}B_\alpha$. If $T=\emptyset$, we are done.\\
Otherwise write the nonempty $B_\alpha$ as $(B_{\alpha_j})_{j\geq1}$, with repetitions if needed, and enumerate $B_{\alpha_j}$ by $(b_{jk})_{k\in\N}$.\\
A diagonal scan of $(b_{jk})$ has image $T$, so Proposition \ref{prop:subsets-and-images-at-most-countable} makes $T$ at most countable.\\
\hlt{Structural Commentary}: Finite pieces cause no difficulty; they can be embedded into countable rows.\\
\hlt{Key Insight}: At most countable indexing $+$ at most countable indexed sets gives at most countable total set.
\end{proof}
\end{corollary}

\begin{theorem} \label{thm:finite-tuples-over-countable-set-countable} \hlt{Finite Tuples over a Countable Set}\\
Let $A$ be countable. For each fixed $n\in\N$, the set of all $n$-tuples $(a_1,\ldots,a_n)$ with $a_k\in A$ is countable.
\begin{proof} Let $B_n\defeq \{(a_1,\ldots,a_n): a_k\in A\}$, $B_1=A$. If $B_{n-1}$ is countable, $B_n=\bigcup\limits_{b\in B_{n-1}}\{(b,a):a\in A\}$, each set $\{(b,a):a\in A\}$ is equivalent to $A$. Hence $B_n$ is countable by Theorem \ref{thm:countable-union-of-countable-sets}.\\
\hlt{Structural Commentary}: Build an $n$-tuple by appending one coordinate to an $(n-1)$-tuple.\\
\hlt{Key Insight}: Induction plus countable unions controls finite products.
\end{proof}
\end{theorem}

\begin{figure}[H]
  \centering
  \begin{tikzpicture}[
  x=0.82cm,
  y=0.82cm,
  >=Stealth,
  every node/.style={font=\small},
  point/.style={circle, fill=black, inner sep=1.25pt},
  active/.style={circle, draw=orange!60!black, fill=orange!12, line width=0.55pt, inner sep=2.2pt},
  guide/.style={orange!75!black, line width=0.65pt},
  scan/.style={->, blue!70!black, line width=0.7pt},
  output/.style={draw=gray!55, fill=gray!5, rounded corners=1pt, inner xsep=5pt, inner ysep=4pt}
]
  \foreach \i in {1,...,5} {
    \node[anchor=east] at (0,-0.42*\i) {$\i$};
    \node[anchor=south] at (0.47*\i,-0.1) {$\i$};
  }
  \node[anchor=south] at (2.65,-0.1) {$\cdots$};

  \draw[guide] (0.94,-0.42) -- (0.47,-0.84);
  \draw[guide] (1.41,-0.42) -- (0.47,-1.26);
  \draw[guide] (1.88,-0.42) -- (0.47,-1.68);

  \foreach \row in {1,...,5} {
    \foreach \col in {1,...,5} {
      \node[point] at (0.47*\col,-0.42*\row) {};
    }
  }

  \foreach \row/\col in {1/1,1/2,2/1,1/3,2/2,3/1,1/4,2/3,3/2,4/1} {
    \node[active] at (0.47*\col,-0.42*\row) {};
  }

  \node[output, anchor=west] (list) at (3.10,-1.05)
    {$(1,1),\,(1,2),\,(2,1),\,(1,3),\,(2,2),\,(3,1),\,\cdots$};
  \draw[scan] (2.70,-1.05) -- (list.west);
\end{tikzpicture}

  \caption{Diagonal enumeration of $\N\times\N$ by finite $i+j$ diagonals.}
  \label{fig:countable-diagonal-enumeration}
\end{figure}

\begin{corollary} \label{cor:countable-cartesian-products} \hlt{Countable Cartesian Products}\\
Let $A,B$ be sets.
\begin{enumerate}[label=(\roman*)]
\setlength{\itemsep}{0pt}
\item If $A$ and $B$ are at most countable, then $A\times B$ is at most countable.
\item If $A$ and $B$ are countable and nonempty, then $A\times B$ is countable.
\end{enumerate}
In particular, $\N\times\N$ is countable.
\begin{proof}
\begin{enumerate}[label=(\roman*)]
\setlength{\itemsep}{0pt}
\item For $a\in A$, set $S_a\defeq\{a\}\times B$. The map $B\to S_a$, $b\mapsto(a,b)$, is bijective, so each $S_a$ is at most countable. Since $A\times B=\bigcup\limits_{a\in A}S_a$ and $A$ is at most countable, Corollary \ref{cor:at-most-countable-union-at-most-countable} gives that $A\times B$ is at most countable.
\item By (i), $A\times B$ is at most countable. Choose $b_0\in B$. The map $A\to A\times B$, $a\mapsto(a,b_0)$, is injective; since $A$ is infinite, $A\times B$ is infinite. Hence $A\times B$ is countable. Taking $A=B=\N$ gives $\N\times\N$ countable.
\end{enumerate}
\hlt{Structural Commentary}: Products decompose into slices.\\
\hlt{Key Insight}: Count slices, then fix a coordinate.
\end{proof}
\end{corollary}

\begin{corollary} \label{cor:rationals-countable} \hlt{The Rationals Are Countable}\\
The set $\Q$ is countable.
\begin{proof} Let $D\defeq\{(p,q)\in\Z\times\Z:q\neq0\}$ and define $\pi:D\to\Q$ by $\pi(p,q)\defeq p/q$.\\
Since $D\subseteq\Z\times\Z$, Corollary \ref{cor:countable-cartesian-products} and Proposition \ref{prop:subsets-and-images-at-most-countable} make $D$ at most countable.\\
Since $\pi(D)=\Q$, $\Q$ is at most countable. The injection $\Z\to\Q$, $n\mapsto n$, shows $\Q$ is infinite; hence $\Q$ is countable.\\
\hlt{Structural Commentary}: Fractions are images of integer pairs.\\
\hlt{Key Insight}: Count pairs, then quotient.
\end{proof}
\end{corollary}
